\chapter{Algebraic Surfaces}\label{algebraic-surfaces}
\chapterstatus{planned}

After curves, surfaces are the first case where birational geometry is not trivial: blowing up points
changes a surface without changing its function field, and the Enriques--Kodaira classification sorts
the minimal models. The intersection pairing on divisors is the main tool throughout.

\noindent\textbf{Prerequisites.} Chapter~\ref{schemes} (Schemes and Coherent Sheaves), Chapter~\ref{divisors} (Divisors and Line Bundles), Chapter~\ref{curves} (Algebraic Curves and Riemann Surfaces) and Chapter~\ref{intersection-theory} (Algebraic Cycles and Chow Groups).

\noindent\textbf{Planned contents.}
\begin{enumerate}
\item The intersection pairing of divisors on a smooth projective surface; the adjunction formula and
the genus of curves on a surface.
\item The Riemann--Roch theorem for surfaces and Noether's formula; the Hodge index theorem and the
Nakai--Moishezon criterion.
\item Blowing up a point, the exceptional curve, the behaviour of the intersection form; factorisation
of birational morphisms; Castelnuovo's contractibility criterion; minimal models.
\item Rational and ruled surfaces: Hirzebruch surfaces, Castelnuovo's rationality criterion, Noether's
theorem on $\mathbb{P}^2$.
\item Cubic surfaces: the $27$ lines, the cubic surface as the blow-up of $\mathbb{P}^2$ in six points,
del Pezzo surfaces and the root systems $E_6, E_7, E_8$.
\item Kodaira dimension; K3 and Enriques surfaces, abelian and bielliptic surfaces, elliptic surfaces
and Kodaira's canonical bundle formula; surfaces of general type and the Bogomolov--Miyaoka--Yau inequality (statement).
\item The Enriques--Kodaira classification (statement and overview).
\item Birational geometry in higher dimension (survey): toric varieties and fans, the minimal model programme, flips, Fano varieties.
\item Exercises with solutions.
\end{enumerate}
